\documentclass[a4paper,14pt]{extarticle}

\input{C:/preamble.tex}

\begin{document}
	
	

	
	Обозначим его вершины
	\[
	\widehat a_1=(0,0),\qquad
	\widehat a_2=(1,0),\qquad
	\widehat a_3=(0,1).
	\]
	Координаты на нём будем обозначать \(\xi,\eta\), чтобы не смешивать с координатами \(x,y\) нашего треугольника.
	
	Базисные функции на эталоне:
	\[
	\widehat\lambda_1(\xi,\eta)=1-\xi-\eta,\qquad
	\widehat\lambda_2(\xi,\eta)=\xi,\qquad
	\widehat\lambda_3(\xi,\eta)=\eta.
	\]
	Проверка в вершинах:
	\[
	\begin{array}{c|ccc}
		&\widehat\lambda_1&\widehat\lambda_2&\widehat\lambda_3\\ \hline
		(0,0)&1&0&0\\
		(1,0)&0&1&0\\
		(0,1)&0&0&1
	\end{array}
	\]
	Это нужный нам узловой базис.
	
	Для эталонных вершин лучше использовать \(\widehat a_i\), поскольку \(a_i^{(h)}\) у тебя обозначают вершины реального треугольника.
	

	
	Пусть вершины нашего треугольника:
	\[
	a_1^{(h)}=(x_1,y_1),\qquad
	a_2^{(h)}=(x_2,y_2),\qquad
	a_3^{(h)}=(x_3,y_3).
	\]
	Здесь рассматриваем фиксированный \(h\), поэтому у координат вершин дополнительный индекс опущен.
	
	Ищем преобразование вида
	\[
	\begin{pmatrix}x\\y\end{pmatrix}
	=
	\underbrace{\begin{pmatrix}
			\alpha&\beta\\
			\gamma&\delta
	\end{pmatrix}}_{A_h}
	\begin{pmatrix}\xi\\\eta\end{pmatrix}
	+
	\underbrace{\begin{pmatrix}b_x\\b_y\end{pmatrix}}_{b_h}.
	\]
	
	**Умножение матрицы на столбец выполняется по строкам:**
	\[
	\begin{pmatrix}
		\alpha&\beta\\
		\gamma&\delta
	\end{pmatrix}
	\begin{pmatrix}\xi\\\eta\end{pmatrix}
	=
	\begin{pmatrix}
		\alpha\xi+\beta\eta\\
		\gamma\xi+\delta\eta
	\end{pmatrix}.
	\]
	Поэтому преобразование означает два обычных равенства:
	\[
	\begin{cases}
		x=\alpha\xi+\beta\eta+b_x,\\
		y=\gamma\xi+\delta\eta+b_y.
	\end{cases}
	\]
	
	Теперь определим все коэффициенты из желаемого соответствия вершин.
	
	Для \((\xi,\eta)=(0,0)\) хотим получить \((x_1,y_1)\):
	\[
	\begin{pmatrix}x_1\\y_1\end{pmatrix}
	=
	\begin{pmatrix}
		\alpha\cdot0+\beta\cdot0\\
		\gamma\cdot0+\delta\cdot0
	\end{pmatrix}
	+
	\begin{pmatrix}b_x\\b_y\end{pmatrix}.
	\]
	Следовательно,
	\[
	b_x=x_1,\qquad b_y=y_1.
	\]
	
	Для \((\xi,\eta)=(1,0)\) хотим получить \((x_2,y_2)\):
	\[
	\begin{pmatrix}x_2\\y_2\end{pmatrix}
	=
	\begin{pmatrix}\alpha\\\gamma\end{pmatrix}
	+
	\begin{pmatrix}x_1\\y_1\end{pmatrix}.
	\]
	Отсюда
	\[
	\alpha=x_2-x_1,\qquad \gamma=y_2-y_1.
	\]
	
	Для \((\xi,\eta)=(0,1)\) хотим получить \((x_3,y_3)\):
	\[
	\begin{pmatrix}x_3\\y_3\end{pmatrix}
	=
	\begin{pmatrix}\beta\\\delta\end{pmatrix}
	+
	\begin{pmatrix}x_1\\y_1\end{pmatrix},
	\]
	поэтому
	\[
	\beta=x_3-x_1,\qquad \delta=y_3-y_1.
	\]
	
	Итак,
	\[
	\boxed{
		\begin{pmatrix}x\\y\end{pmatrix}
		=
		\begin{pmatrix}
			x_2-x_1&x_3-x_1\\
			y_2-y_1&y_3-y_1
		\end{pmatrix}
		\begin{pmatrix}\xi\\\eta\end{pmatrix}
		+
		\begin{pmatrix}x_1\\y_1\end{pmatrix}.
	}
	\]
	

	
	То же умножение можно раскрыть **по столбцам**:
	\[
	\begin{pmatrix}
		x_2-x_1&x_3-x_1\\
		y_2-y_1&y_3-y_1
	\end{pmatrix}
	\begin{pmatrix}\xi\\\eta\end{pmatrix}
	=
	\xi\begin{pmatrix}x_2-x_1\\y_2-y_1\end{pmatrix}
	+
	\eta\begin{pmatrix}x_3-x_1\\y_3-y_1\end{pmatrix}.
	\]
	Следовательно,
	\[
	\boxed{
		\begin{pmatrix}x\\y\end{pmatrix}
		=
		\begin{pmatrix}x_1\\y_1\end{pmatrix}
		+\xi\begin{pmatrix}x_2-x_1\\y_2-y_1\end{pmatrix}
		+\eta\begin{pmatrix}x_3-x_1\\y_3-y_1\end{pmatrix}.
	}
	\]
	
	Смысл такой: **начинаем в первой вершине и прибавляем \(\xi\) раз вектор первого ребра и \(\eta\) раз вектор второго ребра**.
	
	Столбцы \(A_h\) — это как раз два вектора рёбер, выходящих из первой вершины. Они могут иметь другую длину и угол между собой, чем оси эталонного треугольника.
	
	Перегруппировав слагаемые, получаем
	\[
	(x,y)=(1-\xi-\eta)a_1^{(h)}
	+\xi a_2^{(h)}
	+\eta a_3^{(h)}.
	\]
	У точек эталонного треугольника
	\[
	\xi\ge0,\qquad \eta\ge0,\qquad \xi+\eta\le1.
	\]
	Поэтому все три веса неотрицательны и в сумме дают единицу: преобразование действительно переводит весь эталонный треугольник в наш.
	

	
	Наша задача — вычислить, например,
	\[
	\lambda_{p,2}^{(h)}(x,y)
	\]
	в заданной точке **реального** треугольника.
	
	Но известная эталонная функция принимает аргументы \(\xi,\eta\):
	\[
	\widehat\lambda_2(\xi,\eta)=\xi.
	\]
	Значит, сначала нужно ответить: **какая точка эталона соответствует заданной точке \((x,y)\)?**
	
	Для этого решаем полученную систему относительно \(\xi,\eta\):
	\[
	\begin{pmatrix}x-x_1\\y-y_1\end{pmatrix}
	=
	A_h\begin{pmatrix}\xi\\\eta\end{pmatrix},
	\]
	откуда
	\[
	\begin{pmatrix}\xi\\\eta\end{pmatrix}
	=
	A_h^{-1}
	\begin{pmatrix}x-x_1\\y-y_1\end{pmatrix}.
	\]
	
	**5. Раскроем обратную матрицу и умножение**
	
	Для матрицы размера \(2\times2\)
	\[
	\begin{pmatrix}a&b\\c&d\end{pmatrix}^{-1}
	=
	\frac1{ad-bc}
	\begin{pmatrix}d&-b\\-c&a\end{pmatrix}.
	\]
	В нашем случае определитель равен
	\[
	D_h=(x_2-x_1)(y_3-y_1)
	-(x_3-x_1)(y_2-y_1).
	\]
	Для невырожденного треугольника \(D_h\ne0\). Геометрически \(|D_h|\) — удвоенная площадь треугольника; знак зависит от порядка вершин. **В формуле обратной матрицы сохраняем именно \(D_h\), со знаком.**
	
	Тогда
	\[
	A_h^{-1}
	=
	\frac1{D_h}
	\begin{pmatrix}
		y_3-y_1&-(x_3-x_1)\\
		-(y_2-y_1)&x_2-x_1
	\end{pmatrix}.
	\]
	Умножаем каждую строку на столбец:
	\[
	\begin{pmatrix}\xi\\\eta\end{pmatrix}
	=
	\frac1{D_h}
	\begin{pmatrix}
		(y_3-y_1)(x-x_1)-(x_3-x_1)(y-y_1)\\
		-(y_2-y_1)(x-x_1)+(x_2-x_1)(y-y_1)
	\end{pmatrix}.
	\]
	Следовательно,
	\[
	\boxed{
		\xi(x,y)=
		\frac{(y_3-y_1)(x-x_1)-(x_3-x_1)(y-y_1)}{D_h},
	}
	\]
	\[
	\boxed{
		\eta(x,y)=
		\frac{-(y_2-y_1)(x-x_1)+(x_2-x_1)(y-y_1)}{D_h}.
	}
	\]
	
	**6. Как из этого получаются базисные функции**
	
	Теперь просто подставляем найденные координаты в эталонные функции:
	\[
	\boxed{
		\begin{aligned}
			\lambda_{p,1}^{(h)}(x,y)&=1-\xi(x,y)-\eta(x,y),\\
			\lambda_{p,2}^{(h)}(x,y)&=\xi(x,y),\\
			\lambda_{p,3}^{(h)}(x,y)&=\eta(x,y).
		\end{aligned}
	}
	\]
	
	Мы **не умножаем значения функций на матрицу**. Матрица преобразует координаты точки, а функция в соответствующей точке получает то же скалярное значение.
	
	Если обозначить прямое преобразование через \(F_h\), эта подстановка кратко записывается:
	\[
	\lambda_{p,i}^{(h)}=\widehat\lambda_i\circ F_h^{-1}.
	\]
	То есть сначала выполняется \(F_h^{-1}\), затем вычисляется \(\widehat\lambda_i\).
	
	Раскрывая скобки, получаем явный вид всех трёх функций:
	\[
	\boxed{
		\lambda_{p,1}^{(h)}(x,y)=
		\frac{x_2y_3-x_3y_2+(y_2-y_3)x+(x_3-x_2)y}{D_h},
	}
	\]
	\[
	\boxed{
		\lambda_{p,2}^{(h)}(x,y)=
		\frac{x_3y_1-x_1y_3+(y_3-y_1)x+(x_1-x_3)y}{D_h},
	}
	\]
	\[
	\boxed{
		\lambda_{p,3}^{(h)}(x,y)=
		\frac{x_1y_2-x_2y_1+(y_1-y_2)x+(x_2-x_1)y}{D_h}.
	}
	\]
	Вот те самые функции вида \(c_{i,1}+c_{i,2}x+c_{i,3}y\): их коэффициенты теперь выражены через вершины треугольника.
	
	Условие Кронекера сохраняется автоматически. Например, вершина \(a_2^{(h)}\) при обратном преобразовании переходит в \((1,0)\), поэтому значения трёх функций в ней равны \(0,1,0\).
	
	**7. Численный пример всей процедуры**
	
	Возьмём треугольник
	\[
	a_1=(1,1),\qquad a_2=(3,1),\qquad a_3=(2,4).
	\]
	Тогда
	\[
	A=
	\begin{pmatrix}2&1\\0&3\end{pmatrix},
	\qquad
	b=\begin{pmatrix}1\\1\end{pmatrix}.
	\]
	Прямое преобразование:
	\[
	\begin{pmatrix}x\\y\end{pmatrix}
	=
	\begin{pmatrix}2&1\\0&3\end{pmatrix}
	\begin{pmatrix}\xi\\\eta\end{pmatrix}
	+\begin{pmatrix}1\\1\end{pmatrix}
	=
	\begin{pmatrix}2\xi+\eta+1\\3\eta+1\end{pmatrix}.
	\]
	Поэтому
	\[
	y=3\eta+1
	\quad\Longrightarrow\quad
	\eta=\frac{y-1}{3},
	\]
	\[
	x=2\xi+\eta+1
	\quad\Longrightarrow\quad
	\xi=\frac{x-1}{2}-\frac{y-1}{6}.
	\]
	Подставляем:
	\[
	\lambda_1(x,y)
	=1-\frac{x-1}{2}-\frac{y-1}{6},
	\]
	\[
	\lambda_2(x,y)
	=\frac{x-1}{2}-\frac{y-1}{6},
	\qquad
	\lambda_3(x,y)=\frac{y-1}{3}.
	\]
	
	Например, в точке \((2,2)\):
	\[
	\xi=\frac12-\frac16=\frac13,\qquad
	\eta=\frac13.
	\]
	Следовательно,
	\[
	\lambda_1(2,2)=\lambda_2(2,2)=\lambda_3(2,2)=\frac13,
	\]
	и давление в ней выражается как
	\[
	p(2,2)=\frac13\omega_{p,1}
	+\frac13\omega_{p,2}
	+\frac13\omega_{p,3}.
	\]
	
	Таким образом, для каждой точки реального треугольника мы находим её координаты на эталоне и используем их в уже известных формулах \(1-\xi-\eta,\xi,\eta\). **Это позволяет один раз построить базис на эталоне, а затем получать его на каждом треугольнике сетки через координаты вершин.**
	
\end{document}